\section{Mechanism of Forking}
\label{sec:mechanism}

Section~\ref{sec:empirical} left two questions: why the gap follows a power law in context frequency, and why replay keeps widening it. We answer the first with continuation \emph{coverage} and the second with \emph{crowding-out} inside context subspaces, then show that the $n$-gram module is what makes those subspaces independent. Additional experimental details are in Appendix~\ref{app:mechanistic}.

\subsection{Coverage and the frequency law}
\label{sec:coverage}
Section~\ref{sec:forking-factors} found that low-frequency contexts carry most of the gap. Overfitting, in general, hurts on the validation samples that training never showed; in next-token prediction these are the positions whose context occurs in training but whose true continuation never followed it there---\emph{seen context, novel continuation}. Every validation position therefore falls into one of two classes, and the gap splits accordingly. Averaged over the contexts that occur $F$ times in training,
\begin{equation}
\gap(F)=f(F)\,\underbrace{\big[\ell^{\mathrm{val}}_{\mathrm{novel}}(F)-\ell^{\mathrm{tr}}(F)\big]}_{D(F)}
\;+\;\big(1-f(F)\big)\big[\ell^{\mathrm{val}}_{\mathrm{seen}}(F)-\ell^{\mathrm{tr}}(F)\big],
\label{eq:gap-decomp}
\end{equation}
where $f(F)$ is the fraction of validation positions at those contexts whose continuation is novel, $\ell^{\mathrm{tr}}$ is the training loss at the same contexts, and $D(F)$ is the depth to which novel continuations sit below the training fit. The second bracket compares validation and training positions that share both the context row and the continuation, and it is small: novel-continuation positions make up 31\% of the validation stream when contexts are counted as bigrams, yet they carry 73\% of the both-table model's excess validation loss over the no-gram control at step 1000; adding the seen continuations of rare contexts ($F\le 31$) brings this to 77\%, and at trigram granularity novel continuations account for essentially all of it (Appendix~\ref{app:novel-share}). The first bracket carries the frequency law. Replay reinforces the continuations a context was seen with but supplies no evidence about the others, so $D$ grows with the number of passes $j$ (Section~\ref{sec:sharpening}); if it does not otherwise depend on $F$,
\begin{equation}
\gap_j(F)\approx D_j\,f(F).
\label{eq:gap_j}
\end{equation}
The measured gaps of Figure~\ref{fig:frequency} fit
\begin{equation}
\gap^{\mathrm{bi}}(F)\approx 7.21\,F^{-0.233},\qquad
\gap^{\mathrm{tri}}(F)\approx 3.76\,F^{-0.318}.
\label{eq:gap-fits}
\end{equation}
Natural language supplies the power law in $f$. Good--Turing estimation \citep{good1953population,gale1995good} gives the probability that a context seen $F$ times is followed by an unseen continuation as $\hat f(F)=N_1(F)/F$, where $N_1$ counts continuation types that followed it exactly once---a statistic of the text, not of any model. If conditional continuation frequencies have a Zipf-like tail $p_{(r)}\propto r^{-s}$, this novel-continuation rate decays as $F^{-\alpha}$ with $\alpha=1-1/s$ \citep{gnedin2007occupancy}; the derivation is in Appendix~\ref{app:math}. On our corpus the estimate, averaged over the contexts at each count $F$, is a clean power law over the window $F\le10^3$ in which the gap is measured (Appendix Figure~\ref{fig:b_t_MF}),
\begin{equation}
\hat f^{\mathrm{bi}}(F)\approx 0.95\,F^{-0.287},\qquad
\hat f^{\mathrm{tri}}(F)\approx 0.97\,F^{-0.314}.
\label{eq:gt-fits}
\end{equation}
The exponents read off corpus counts thus land within $0.05$ (bigram) and $0.01$ (trigram) of those the model produces, and a single depth $D$ per branch turns $\hat f(F)$ into the measured gap: across the seven observed frequency bins the ratio $\operatorname{gap}/\hat f$ is constant to within 9\% and 18\% (Appendix Figure~\ref{fig:gt-kernel}). For bigram contexts that ratio is about 8 nats---the size of the seen-to-novel log-odds shift visible on a single context in Figure~\ref{fig:walk-their-dog}. The frequency law of Section~\ref{sec:forking-factors} is therefore a property of the corpus: rare contexts have more of their continuation mass unseen. Coverage does not account for everything---about a quarter of the excess loss falls on seen continuations of frequent contexts, and for trigram contexts $D$ drifts slightly with $F$ (Appendix~\ref{app:gt-fits})---but it fixes the exponent.

\begin{figure}[!t]
\centering
\begin{subfigure}[b]{0.58\linewidth}
\centering
\includegraphics[width=\linewidth]{figures/walk_their_dog_loss.pdf}
\subcaption{}
\label{fig:walk-their-dog}
\end{subfigure}\hfill
\begin{subfigure}[b]{0.38\linewidth}
\centering
\includegraphics[width=\linewidth]{figures/fig09_crowd_ce_seen_novel.pdf}
\subcaption{}
\label{fig:crowd-ce}
\end{subfigure}
\caption{Crowding-out: (a) context-only continuation losses for ``walk their dog''; (b) validation cross-entropy of train-seen and novel continuations, pooled over contexts.}
\label{fig:crowding}
\end{figure}

\subsection{Crowding-out under replay}
\label{sec:sharpening}
Coverage says what is missing. The more important questions are why replay turns the missing coverage into loss, why the $n$-gram module is needed for it, and how the damage evolves with the number of passes. Our answer is that forking is ordinary overfitting confined to very small problems. For a context $c_i$, call the next-token logits computed from it its \emph{context subspace}, and split the vocabulary into the continuations of $c_i$ seen in training, those that appear only in validation, and the rest (Appendix~\ref{app:crowding-derivation}). If training shows a single continuation $C$, the loss in that subspace is $-\ln p_C$; a gradient step raises the logit of $C$ by $\eta(1-p_C)$ and, through the softmax, lowers every other logit by $\eta p_v$. Mass flows from the unseen continuations onto the seen one, so the validation loss on plausible-but-unseen continuations rises by the same shift that lowers the training loss, and the seen-to-unseen odds grow by a factor above one at every pass with nothing to stop them while $p_C<1$ (Appendix Equation~\eqref{eq:zt}). We call this \emph{crowding-out}; it is the growth of the depth $D_j$ in Equation~\eqref{eq:gap-decomp} with replay.

Figure~\ref{fig:walk-their-dog} shows the picture on one context. ``walk their dog'' is followed by \emph{down} once in training and, in the fixed validation window, once each by \emph{daily} and \emph{for}. Probed with context-only input, the probability of \emph{down} rises by more than an order of magnitude over five passes while \emph{daily} and \emph{for} are compressed---\emph{daily} by orders of magnitude---and so is the total mass outside the continuations observed in either sample. Figure~\ref{fig:crowd-ce} shows that the context is typical. Pooling validation positions whose bigram context occurs in training, split by whether the continuation was seen there (405{,}058 positions) or not (158{,}790), both models do well on seen continuations: the control's cross-entropy falls from 4.3 to 2.8 nats and the both-table model's stays near 3.4. On novel continuations they diverge: the control still improves, 6.9 to 5.6 nats, while the both-table model deteriorates from 7.5 to 9.0 with a step at every epoch boundary---the staircase of Figure~\ref{fig:staircase}. Consistently, at epoch~3 the both-table model's predictive entropy on validation is 2.93 nats against a cross-entropy of 5.26, whereas the control's are 3.59 and 3.58 (same runs as Figure~\ref{fig:crowd-ce}): far more confident than accurate, the signature of a model that has committed to few continuations per context \citep{carlsson2025hyperfitting}. Crowding-out also explains why freezing the tables does not stop forking (Section~\ref{sec:interventions}): the subspace coordinate is already in place, and the backbone can continue the descent on its own. Coverage thus fixes the size of the target, $f\propto F^{-\alpha}$, crowding-out fixes its depth, $D_j$, and their product is the gap.

\subsection{Context subspaces are independent}
\label{sec:context-locality}
\begin{figure}[!t]
\centering
\begin{subfigure}[b]{0.45\linewidth}
\includegraphics[width=\linewidth]{figures/fig07a_context_mask.pdf}
\subcaption{Gap contribution.}
\label{fig:ctxmask-gap}
\end{subfigure}\hfill
\begin{subfigure}[b]{0.45\linewidth}
\includegraphics[width=\linewidth]{figures/fig07b_seen_mass.pdf}
\subcaption{Mass on train-seen continuations.}
\label{fig:ctxmask-mass}
\end{subfigure}
\caption{Random context loss masking.}
\label{fig:context-subspace-mask}
\end{figure}
Section~\ref{sec:sharpening} treats each context subspace as evolving on its own, with the sharpening within it as the engine of forking. Both statements are claims about what the $n$-gram module buys: a private coordinate per context, so that updates driven by one context barely touch the conditional distribution of another. We test this by ablation. We hold out a random 2\% of training trigram contexts and exclude the next-token loss at their positions from the start of epoch~2 or epoch~3 onward, keeping inputs, trainable parameters and the loss denominator unchanged; a disjoint 2\% serves as unmasked peers, and each masked run is paired with normal training from the same checkpoint. Figure~\ref{fig:ctxmask-gap} plots their gap contribution $C_S=\sum_{i\in S}q_i[\bar\ell_{\mathrm{val},i}-\bar\ell_{\mathrm{tr},i}]$ on fixed batches. Until the mask fires the masked curves track normal training; afterwards they freeze---for two and three further passes---while the unmasked and normal curves take a fresh jump at every replay boundary. The gap of a masked context stops growing the moment its own replay stops, although the remaining 98\% of contexts keep training throughout: gap growth is generated inside each subspace, not imported from the others.

Figure~\ref{fig:ctxmask-mass} follows the sharpening coordinate itself, the predictive mass on the continuations of each context that training has seen. In the n-gram model the unmasked contexts pour mass into their seen continuations pass after pass ($0.43\to0.69$ by epoch~4) while the masked ones freeze ($0.43\to0.40$): once a context's own loss is removed, nothing else sharpens its subspace. Without the $n$-gram module the mask changes nothing ($0.22\to0.39$ masked vs.\ $0.41$ unmasked), and replay raises the mass on seen and validation-only continuations alike ($0.013\to0.06$): every context's conditional distribution lives in parameters shared with all others, and repetition improves it holistically. The table therefore partitions optimization into weakly interacting context subspaces; inside a subspace, replay-driven sharpening proceeds unchecked by the continuations of other contexts, which is exactly the over-memorization assumed in Section~\ref{sec:sharpening}. Masked contexts are not immutable---the shared backbone still drifts---but that drift moves train and validation losses together and never enters the gap of Figure~\ref{fig:ctxmask-gap}.

